\documentclass[11pt,a4paper]{article}

\usepackage[T1]{fontenc}
\usepackage[utf8]{inputenc}
\usepackage[ngerman]{babel}
\usepackage{amsmath,amssymb,amsthm,amsfonts}
\usepackage{mathtools}
\usepackage{graphicx}
\usepackage{hyperref}
\usepackage{geometry}
\usepackage{enumitem}
\usepackage{booktabs}
\usepackage{physics}

\geometry{margin=2.5cm}

\title{
	Vom EABC-Tetraeder zur modularen Spin-Geometrie\\
	Eine Übersicht über das Bamberger Modell,
	goldene Verfeinerungen und arithmetische Drehimpulse
}

\author{
	Thomas Hoffbauer\\
	Bamberg, Deutschland
}

\date{\today}

\begin{document}
	
	\maketitle
	
	\begin{abstract}
		Dieser Übersichtsartikel fasst die gegenwärtige Struktur des Bamberger Modells zusammen.
		Ausgehend von der EABC-Klassifikation der natürlichen Zahlen wird eine geometrische
		Interpretation entwickelt, in der die vier EABC-Zustände die Ecken eines regulären
		Tetraeders bilden.
		
		Der Tetraeder wird als diskreter Spinraum interpretiert.
		Die natürliche Erweiterung über den Goldenen Schnitt führt zu einer ikosaedrischen
		Verfeinerung, welche erstmals geschlossene Umläufe und damit einen arithmetischen
		Bahndrehimpuls ermöglicht.
		
		Weiterhin wird gezeigt, wie die Dedekindsche Eta-Funktion,
		die Zahl 24 der Hurwitz-Einheiten, Ramanujans Identitäten
		und die Ikosaedersymmetrie der Gruppe $A_5$
		in einem gemeinsamen Rahmen erscheinen.
		
		Das resultierende Modell trennt die Struktur natürlicher Zahlen in
		Skala, Spin, Bahnimpuls und Spinorphase und liefert damit einen
		möglichen Ausgangspunkt für eine diskrete arithmetische Geometrie.
	\end{abstract}
	
	\tableofcontents
	
	\section{Einleitung}
	
	Das Bamberger Modell entstand ursprünglich als Untersuchung der
	Restklassenstruktur von Primzahlen modulo 12.
	
	Die vier Einheiten
	
	\[
	(\mathbb Z/12\mathbb Z)^\times
	=
	\{1,5,7,11\}
	\]
	
	werden durch die Symbole
	
	\[
	E,\quad A,\quad B,\quad C
	\]
	
	repräsentiert.
	
	Im Verlauf der Untersuchungen zeigte sich,
	dass diese vier Zustände nicht lediglich eine
	Restklassenklassifikation darstellen,
	sondern eine überraschend reichhaltige geometrische Struktur besitzen.
	
	Der vorliegende Artikel fasst die daraus entwickelte
	Spin- und Bahnimpulsinterpretation zusammen.
	
	\section{Die EABC-Algebra}
	
	Die grundlegende Multiplikation lautet
	
	\[
	A^2=B^2=C^2=E
	\]
	
	sowie
	
	\[
	AB=C,
	\qquad
	BC=A,
	\qquad
	CA=B.
	\]
	
	Diese Struktur ist isomorph zur Klein-Vierergruppe
	
	\[
	V_4.
	\]
	
	Die Algebra ist kommutativ und beschreibt die elementaren
	Orientierungszustände des Modells.
	
	\section{Der EABC-Spinraum}
	
	Wir definieren
	
	\begin{align}
		E &= (1,1,1),\\
		A &= (-1,-1,1),\\
		B &= (1,-1,-1),\\
		C &= (-1,1,-1).
	\end{align}
	
	\subsection{Tetraederstruktur}
	
	Zwischen allen Punktpaaren gilt
	
	\[
	d(X,Y)=2\sqrt2.
	\]
	
	Damit besitzen sämtliche Kanten dieselbe Länge.
	
	\begin{section}
		Die vier EABC-Zustände bilden die Ecken eines regulären Tetraeders.
	\end{section}
	
	\begin{proof}
		Alle sechs Kantenlängen stimmen überein.
	\end{proof}
	
	\subsection{Spinneutralität}
	
	Es gilt
	
	\[
	E+A+B+C=0.
	\]
	
	Daher verschwindet der Gesamtspin eines vollständigen Vierlings.
	
	\[
	\boxed{
		S_{\text{Vierling}}=0
	}
	\]
	
	Der Primvierling erscheint somit als spinneutrales Singulett.
	
	\section{Spin und Bahndrehimpuls}
	
	\subsection{Spin}
	
	Der Spin wird als Punkt im Tetraeder interpretiert:
	
	\[
	S=s(H).
	\]
	
	Dabei bestimmt der halbkleinsche Anteil $H$
	die Orientierung.
	
	\subsection{Bahndrehimpuls}
	
	Für einen Übergang
	
	\[
	X\rightarrow Y
	\]
	
	definieren wir
	
	\[
	L_{XY}
	=
	s(X)\times s(Y).
	\]
	
	Für einen geschlossenen Umlauf
	
	\[
	P_1\rightarrow P_2\rightarrow \cdots
	\rightarrow P_n\rightarrow P_1
	\]
	
	ergibt sich
	
	\[
	L=
	\sum_i P_i\times P_{i+1}.
	\]
	
	Somit entsteht Bahndrehimpuls ausschließlich durch Bewegung
	auf der zugrunde liegenden Geometrie.
	
	\section{Geschlossene Umläufe des Tetraeders}
	
	Der vollständige Tetraeder besitzt
	
	\begin{itemize}
		\item vier Dreieckszyklen,
		\item drei Hamilton-Umläufe.
	\end{itemize}
	
	Die Dreieckszyklen reproduzieren erneut tetraedrische Richtungen.
	
	Die Hamilton-Umläufe erzeugen die drei kartesischen Achsen
	
	\[
	\pm x,\quad \pm y,\quad \pm z.
	\]
	
	Dies stellt die erste Erscheinung einer
	Bahndrehimpulsstruktur dar.
	
	\section{Quaternionische Erweiterung}
	
	Wir identifizieren
	
	\[
	E\leftrightarrow 1,
	\qquad
	A\leftrightarrow i,
	\qquad
	B\leftrightarrow j,
	\qquad
	C\leftrightarrow k.
	\]
	
	Damit wird die kommutative Klein-Vierer-Algebra
	zu einer quaternionischen Struktur erweitert.
	
	Nun gilt
	
	\[
	ij=k,
	\qquad
	ji=-k.
	\]
	
	Hier erscheint erstmals eine Orientierung
	der Umlaufrichtung.
	
	\subsection{Spinorphase}
	
	Die Zustände
	
	\[
	\pm E,\pm A,\pm B,\pm C
	\]
	
	bilden die Quaternionengruppe $Q_8$.
	
	Dadurch entsteht eine natürliche Spinorstruktur.
	
	\section{Der Goldene Schnitt}
	
	Der Goldene Schnitt
	
	\[
	\varphi
	=
	\frac{1+\sqrt5}{2}
	\]
	
	taucht im Tetraeder selbst noch nicht auf.
	
	Er erscheint erstmals bei der Verfeinerung
	zum Ikosaeder.
	
	\subsection{Ikosaederkoordinaten}
	
	Die Ecken des Ikosaeders können geschrieben werden als
	
	\[
	(0,\pm1,\pm\varphi),
	\]
	
	\[
	(\pm1,\pm\varphi,0),
	\]
	
	\[
	(\pm\varphi,0,\pm1).
	\]
	
	Die Symmetriegruppe ist
	
	\[
	A_5.
	\]
	
	\section{Der Ikosaeder-Fixpunkt}
	
	Die Verfeinerung vom Tetraeder zum Ikosaeder
	stellt den Übergang von diskreten
	Spinrichtungen zu orbitalen Umläufen dar.
	
	Insbesondere entstehen natürliche
	Pentagon-Umläufe.
	
	Für einen solchen Umlauf gilt
	
	\[
	L_5
	=
	\sum_i P_i\times P_{i+1}.
	\]
	
	Die resultierende Achse verläuft entlang einer
	goldenen Richtung.
	
	\[
	L_5
	\parallel
	(0,1,\varphi).
	\]
	
	\section{Dedekinds Eta-Funktion}
	
	Wir betrachten
	
	\[
	\eta(\tau)
	=
	q^{1/24}
	\prod_{n=1}^{\infty}(1-q^n),
	\qquad
	q=e^{2\pi i\tau}.
	\]
	
	Die Transformation
	
	\[
	\eta(\tau+1)
	=
	e^{\pi i/12}\eta(\tau)
	\]
	
	erzeugt unmittelbar die Zahlen
	
	\[
	12
	\]
	
	und
	
	\[
	24.
	\]
	
	\subsection{Hurwitz-Einheiten}
	
	Die 24 Einheiten des Hurwitz-Gitters
	bilden die natürliche geometrische Entsprechung
	zur 24-fachen Struktur der Eta-Funktion.
	
	\section{Ramanujan und die goldene Resonanz}
	
	Ramanujans Rogers-Ramanujan-Kettenbruch verbindet
	das Nome
	
	\[
	q=e^{-2\pi}
	\]
	
	mit dem Goldenen Schnitt.
	
	Formal:
	
	\[
	R(e^{-2\pi})
	=
	\sqrt5\,\varphi-\varphi.
	\]
	
	Dadurch entsteht eine analytische Verbindung zwischen
	
	\begin{itemize}
		\item modularer Geometrie,
		\item Goldenem Schnitt,
		\item Ikosaedersymmetrie.
	\end{itemize}
	
	\section{Die vollständige Struktur}
	
	Das Modell ordnet jeder natürlichen Zahl
	
	\[
	n
	=
	(G,H,V,\varepsilon)
	\]
	
	vier unabhängige Komponenten zu.
	
	\begin{center}
		\begin{tabular}{ll}
			\toprule
			Komponente & Bedeutung\\
			\midrule
			$G$ & glatte Skala\\
			$H$ & tetraedrischer Spin\\
			$V$ & ikosaedrische Verfeinerung\\
			$\varepsilon$ & Spinorphase\\
			\bottomrule
		\end{tabular}
	\end{center}
	
	Dabei gilt
	
	\[
	S=S(H),
	\]
	
	\[
	L=L(G,V),
	\]
	
	\[
	J=S+L.
	\]
	
	\section{Ausblick}
	
	Die vorliegende Struktur legt nahe, dass die
	arithmetischen Eigenschaften natürlicher Zahlen
	eine diskrete Geometrie tragen, die sich in mehreren
	Stufen organisiert:
	
	\[
	V_4
	\rightarrow
	Q_8
	\rightarrow
	A_4
	\rightarrow
	A_5
	\rightarrow
	\eta(\tau)
	\rightarrow
	SU(2).
	\]
	
	Eine zentrale offene Frage besteht darin,
	ob die kontinuierliche Spin-Geometrie von $SU(2)$
	als Grenzübergang der tetraedrisch-ikosaedrischen
	Verfeinerung verstanden werden kann.
	
	Darüber hinaus bleibt zu untersuchen,
	ob Hecke-Operatoren, modulare Formen
	und die Spektren der Riemannschen Zetafunktion
	natürlich in diese Geometrie eingebettet werden können.
	
	\bibliographystyle{plain}
	
	\section{Fermionische Hebung der Kern-Hülle-Struktur}
	
	Die bisherige EABC-Struktur besitzt auf der sichtbaren Restklassenebene die
	kommutative Klein-Vierer-Algebra
	\[
	A^2=B^2=C^2=E,
	\qquad
	AB=C,\quad BC=A,\quad CA=B.
	\]
	Insbesondere gilt dort
	\[
	BC=CB=A.
	\]
	Diese Ebene ist daher noch nicht fermionisch. Sie beschreibt lediglich die
	sichtbare Halbklein-Klasse eines Zustands.
	
	Die fehlende Orientierung entsteht erst durch die quaternionische Hebung
	\[
	E\mapsto 1,\qquad A\mapsto i,\qquad B\mapsto j,\qquad C\mapsto k.
	\]
	Dann gilt
	\[
	ij=k,\qquad jk=i,\qquad ki=j,
	\]
	aber zugleich
	\[
	ji=-k,\qquad kj=-i,\qquad ik=-j.
	\]
	Damit wird insbesondere
	\[
	BC=+A,
	\qquad
	CB=-A.
	\]
	
	Dies ist der erste fermionische Vorzeichenwechsel des Modells. Die sichtbare
	EABC-Algebra bleibt kommutativ, während ihre quaternionische Hebung eine
	orientierte Spinorstruktur trägt.
	
	Der Hüllendipol wird definiert als
	\[
	D_{BC}=B+C.
	\]
	Mit
	\[
	B=(1,-1,-1),
	\qquad
	C=(-1,1,-1)
	\]
	ergibt sich
	\[
	D_{BC}=(0,0,-2).
	\]
	Die Gegenorientierung ist
	\[
	-D_{BC}=(0,0,2).
	\]
	Somit besitzt der Hüllendipol zwei natürliche Spinrichtungen.
	
	Die Multiplikation der beiden Hüllenkomponenten liefert dagegen den
	Bindungskanal:
	\[
	B\circ C=+A,
	\qquad
	C\circ B=-A.
	\]
	Daher ist \(A\) nicht bloß ein statischer EABC-Zustand, sondern ein
	orientierter Bindungskanal zwischen Kern und Hülle.
	
	Wir schreiben folglich
	\[
	E=\text{Kernkanal},
	\]
	\[
	D_{BC}=B+C=\text{Hüllendipol},
	\]
	\[
	\pm A=\text{orientierte Kern-Hülle-Bindung}.
	\]
	
	Das Vorzeichen wird in der modularen Spinorphase gespeichert. Da die
	Dedekindsche Eta-Funktion die Transformation
	\[
	\eta(\tau+1)=e^{\pi i/12}\eta(\tau)
	\]
	besitzt, erscheint eine natürliche \(\mathbb Z_{24}\)-Phasenstruktur.
	Wir setzen
	\[
	+A\leftrightarrow \varepsilon=0,
	\qquad
	-A\leftrightarrow \varepsilon=12
	\pmod{24}.
	\]
	Die Vertauschung
	\[
	BC\longrightarrow CB
	\]
	entspricht daher einem Halbumlauf in der 24-Phase:
	\[
	\varepsilon\longrightarrow \varepsilon+12.
	\]
	Nach zwei Vertauschungen kehrt die Phase zurück:
	\[
	12+12=24\equiv0.
	\]
	Dies ist die spinorische Signatur
	\[
	360^\circ\mapsto -1,
	\qquad
	720^\circ\mapsto +1.
	\]
	
	Damit erhält das Modell drei deutlich getrennte Ebenen:
	\[
	V_4=\text{sichtbare EABC-Restklassenstruktur},
	\]
	\[
	Q_8=\text{quaternionische Vorzeichenstruktur},
	\]
	\[
	\mathbb Z_{24}=\text{modulare Spinorphase}.
	\]
	
	Die verschärfte Kern-Hülle-Regel lautet somit:
	\[
	\boxed{
		E=\text{Kernkanal},
	}
	\]
	\[
	\boxed{
		B+C=\text{Hüllendipol},
	}
	\]
	\[
	\boxed{
		B\circ C=+A,\qquad C\circ B=-A,
	}
	\]
	\[
	\boxed{
		\pm A=\text{orientierte elektromagnetische Bindung}.
	}
	\]
	
	Die fermionische Struktur des Modells liegt also nicht bereits in der
	kommutativen EABC-Algebra, sondern in ihrer quaternionischen Hebung. Die
	Halbklein-Klasse beschreibt den sichtbaren Zustand; die Quaternionen liefern
	das Vorzeichen; die Eta-\(24\)-Phase speichert die Spinorphase.
	
	\begin{thebibliography}{99}
		
		\bibitem{ramanujan}
		S. Ramanujan,
		\emph{Collected Papers},
		Chelsea Publishing.
		
		\bibitem{klein}
		F. Klein,
		\emph{Vorlesungen über das Ikosaeder},
		Teubner.
		
		\bibitem{hurwitz}
		A. Hurwitz,
		\emph{Über die Komposition der quadratischen Formen},
		Mathematische Annalen.
		
		\bibitem{conway}
		J. H. Conway, D. A. Smith,
		\emph{On Quaternions and Octonions},
		A K Peters.
		
	\end{thebibliography}
	
\end{document}